\documentclass[11pt]{article}
\usepackage[margin=1in]{geometry}
\usepackage{amsmath,amssymb,amsthm,hyperref,xurl}
\newtheorem{proposition}{Proposition}
\newtheorem{lemma}{Lemma}
\newcommand{\PA}{\mathsf{PA}}
\newcommand{\I}{\mathsf I}
\newcommand{\M}{\mathsf M}
\newcommand{\U}{\mathsf U}
\title{A guard-free separation gate: soundness and the remaining completeness obligation}
\date{13 November 2026 (Pacific time)}
\begin{document}\maketitle
\section{Precise proposal}
Put $p(a,b)=(a+b)(a+b+1)/2+b$, $\I(a,b)=4p(a,b)$,
$\M(a,b)=4p(a,b)+1$, and $\U(a,B)=4p(B,a)+2$.
For arbitrary numerical inputs $a,b,A$, let
\[
 z=f(a,b,A)=\M(a,\I(b,A)).
\]
The property formula still has code $A$; only the separator variable label
changes from $A$ to $z$. Define
\begin{align*}
 P&=A\to\M(z,b),\\
 S&=\U\bigl(b,(\M(a,b)\to P)\wedge(\M(z,b)\to A)\bigr),\\
 Q&=\U(z,S\to\neg P),\\
 V&=\M(a,b)\to\U(z,\M(a,b)),\\
 G_f(a,b,A)&=\neg(V\to Q).
\end{align*}
Classical abbreviations here retain the numerical false code $1=\M(0,0)$.
This gate has no input guard. The original, separately emitted
$\Gamma(A)=A\to\U(A,A)$ can be retained unchanged.

\begin{lemma}
$f(a,b,A)>a,b,A$ for all natural inputs, including $(0,0,0)$.
The assertions are provable in PA.
\end{lemma}
\begin{proof}
Cantor pairing satisfies $p(u,v)\ge u,v$. Consequently
\[
 f(a,b,A)\ge4a+1,
 \qquad f(a,b,A)\ge4\I(b,A)+1=16p(b,A)+1
 \ge16b+1,16A+1.
\]
At the all-zero input, $f(0,0,0)=1$; there is no code-zero collapse.
Every variable label in the decoded formula $A$ is strictly less than $A$,
so the variable with label $z$ is absent from that formula.
\end{proof}

\section{The soundness direction is available}
Use the inherited PA translation: membership $\M(a,b)$ becomes the binary
digit predicate $B(x_a,x_b)$; equality becomes arithmetic equality; the
formula constructors become implication and quantification, with the
usual special proof for $\I(0,0)=0$.

\begin{proposition}
There is a primitive-recursive supplier of PA proofs of the translation
of every gate $G_f(a,b,A)$. PA proves the supplier's syntactic correctness.
Thus the modified numerical calculus, with the other supplied heads and
rules unchanged, has a PA-verifiable refutation translation into PA.
\end{proposition}
\begin{proof}
Write $\phi=T(A)$. The strict bounds imply $z\ne a,b$ and that $x_z$
is not free in $\phi$. The translated $V$ is an ordinary first-order
vacuity theorem. To show $\neg T(Q)$, it suffices to prove
$\exists x_z(T(S)\wedge T(P))$.

If the labels $a,b$ differ, apply the inherited finite bit-comprehension
supplier with bound $x_a+x_b+1$ and matrix
\[
 (B(x_a,y)\lor y=x_b)\wedge\phi[y/x_b].
\]
Its coded set $s$ satisfies the two implications in $T(S)$ and the
distinguished-point condition $T(P)$. Bounded and unbounded versions agree
because a set bit of $x_a$ has index less than $x_a$, and $s<2^n$ has no set
bit at an index at least $n$.

If the labels coincide, the first antecedent under the inner binder is
$B(y,y)$, which PA refutes. Use bound $x_b+1$ and matrix
$(y=x_b)\wedge\phi[y/x_b]$ instead. This supplies the required remaining
implication and point. All substitutions are capture avoiding; the original
free value of $x_b$ is preserved as a parameter in the proof template.

In either case choose $x_z=s$. Freshness ensures this choice does not
change $x_a,x_b$, or $\phi$. Then classical reasoning derives $T(G_f)$.
The construction is the existing PA finite-comprehension proof template,
now with a separately computed fresh witness label. It uses one PA
induction instance on the finite-set bound, with the induction formula
constructed from the input property; it uses no truth predicate or reflection.
Composing these suppliers through a finite derivation gives the claimed
uniform refutation translation.
\end{proof}

\section{Completeness is not established by these bounds}
The old numerical gate directly supplies
\[
 V(a,b,c)=\M(a,b)\to\U(c,\M(a,b))
 \quad(c\ne a,b)
\]
for every permitted label $c$. The modified gate supplies this only when
$c=f(a,b,A)$, hence only for labels larger than $a,b$. For example,
$V(1,2,0)$ has no direct instance in the modified family.
This absence of a direct instance is not a nonderivability theorem.

The supplied PA-refutation compiler uses arbitrary-label membership vacuity
inside negative-property specialization, reverse membership rewriting,
vacuity through Boolean and quantified formulas, and binder transport.
Therefore it cannot simply be imported unchanged. An unrestricted
alpha-conversion or first-order completeness theorem cannot be assumed
while deriving the restricted numerical calculus's logical interfaces.

Two legitimate routes remain: derive every required arbitrary-label
vacuity instance in the new calculus, or provide a new PA-refutation
compiler whose bound-variable discipline uses only available gate instances.
Neither route is proved in this note. Hence the guard-free gate is a sound
research proposal, not a completed Con(PA) equivalence or an accepted
smaller consistency machine.

\section{Prior-art comparison}
The proposed separator/property split is distinct from the following
recorded attempts, but the high-label vacuity obstruction is inherited.
All paths refer to ProveIt commit
\texttt{317b07268681b44ff352d173a162e2a8c241b2fe}, below
\path{Computability/BusyBeaver/Research/PA/docs/}.
\begin{itemize}
\item \path{pa_consistency_cross_archive/theory_results.tex} pads the
property to $A'=\I(\I(D+1,\I(H,D+1)),A)$, proving $A'>D,H$ and propositional
equivalence to $A$. It explicitly leaves arbitrary-vacuity recovery open;
its 210-control implementation attempt fails a source contract at site 2413.
\item \path{pa_tm_extended_research_2026-11-13/report/pa_tm_extended_research.tex},
``Why a larger fresh label does not simply remove the guard,'' considers
$C=c+H+D+1$ and explains the same missing low-label vacuity interface.
\item \path{pa194c_conpa/report/pa194c_conpa.tex} lists a cheaply constructed
automatically fresh label as a further research question, with repeated
label construction and storage identified as physical-cost concerns.
\item The concrete dependency on unrestricted fresh labels is visible in
\path{../jamesyc-upstream/pa196/PA196/Logic/Reverse.lean}, especially
\texttt{negativeSpecialization} and \texttt{chiVacuity}, and in the ordinary
diagonal-calibration replacement argument of
\path{pa195_conpa_2026-11-13/report/sections/pa_calibration.tex}.
\end{itemize}
\end{document}
